\begin{table}[tbp]
\caption{The book's cumulative contribution and the questions returned to theory. This is an analytical synthesis of the preceding chapters, not an additional empirical study.}
\label{tab:ch11-evidence}
\centering\small\renewcommand{\arraystretch}{1.18}
\begin{tabularx}{\linewidth}{@{}>{\raggedright\arraybackslash}p{.19\linewidth}>{\raggedright\arraybackslash}X>{\raggedright\arraybackslash}X@{}}
\toprule
\textbf{Evidence base} & \textbf{Established contribution} & \textbf{Implication for the framework}\\\midrule
Chapters 1--3: foundations & Objects, value, reference conditions, and an evaluation-system design. & A measure becomes evaluative evidence through a question and an intended use.\\[3pt]
Chapter 4: Activity & A specified representation of recorded collaborative participation. & Event selection and aggregation define the observed dimension; they do not exhaust contribution.\\[3pt]
Chapter 5: OpenRank & Relational evaluation with explicit graph and algorithmic choices. & Influence depends on the represented relations and comparison scope.\\[3pt]
Chapter 6: OpenDigger & Recurrent computation, data services, and a contribution-recognition deployment. & Infrastructure and participant responses belong in the account of evaluation.\\[3pt]
Chapter 7: recognition & Educational applications combining relational and contextual evidence. & Contribution, achievement, and learning require distinct warrants.\\[3pt]
Chapter 8: sustainability & Evidence about archival, maintenance, exit, transfer, and renewal. & The continuity of collective capability can exceed the life of a repository.\\[3pt]
Chapter 9: governance & Analysis of transitions, organizational contributions, and attribution. & Object boundaries and institutional conditions may change during evaluation.\\[3pt]
Chapter 10: think tanks & Reports, observatories, visual investigation, and knowledge products. & Evidence must retain its meaning as it enters institutional use.\\
\bottomrule
\end{tabularx}
\end{table}
